% Glossary terms from ch24: lines of the form \item[Term] Definition.
\item[Hybrid collision-avoidance architecture] The four-layer design of the drone system: a global swarm planner (\cbs/\ecbs) for the nominal time-indexed paths, a prediction layer (Kalman filter/LSTM) for unknown objects, a local safety layer (\orca/DWA) for short avoidance maneuvers, and a replanning layer (\dstarlite/\astar) from the current state, joined by a per-drone executive.
\item[Executive] The per-drone state machine that runs every control cycle and decides which layer is in charge: it follows the plan, hands the velocity to the local layer when a predicted conflict is inside the safety horizon, reconnects to the plan when the conflict has cleared, replans when reconnection fails, and re-checks conflicts after every change of a path.
\item[Hybrid world model] The shared representation of the layers: a grid with discrete time (paths with their start times, the reservation table, the blocked layers) and continuous positions, velocities and predictions, with one clock and one frame and explicit conversions between the two sides.
\item[Reference position $\pos^{\mathrm{ref}}_i(t)$] The continuous position that a time-indexed path prescribes at clock time $t$, obtained by linear interpolation between the centers of consecutive cells; the preferred velocity of the local layer is a pursuit of the reference a short lookahead ahead.
\item[Drift $d_i(t)$] The distance between a drone and its reference position; a bound $d_{\max}$ on the drift is one of the two exits from the Avoiding state into Replanning.
\item[Safety horizon $\tau_h$] The look-ahead window of the trigger: a predicted conflict is inside the horizon when the predicted separation minus the uncertainty inflation drops below the safety radius within $\tau_h$ seconds; chosen at least as long as the \orca horizon and the time the drone needs to execute a velocity change.
\item[Predicted time to collision $t_c$] The first look-ahead at which the margin $\norm{\tilde{\pos}_i(s) - \hat{\pos}_{T+h}} - \kappa_p\sigma(h) - R_{\mathrm{safe}}$ is negative; the trigger fires when $t_c \le \tau_h$.
\item[Inflated safety radius $d^{\mathrm{safe}}_h$] The required separation $R_{\mathrm{safe}} = r_A + r_B + \Delta r$ plus $\kappa_p\sqrt{\lambda_{\max}(\mat{\Sigma}_h)}$, the radius of the disk that contains the $p$-confidence ellipse of the prediction at step $h$; keeping this distance from the predicted mean is a per-step chance constraint with probability $p$.
\item[Blocked layer $\mathcal{B}_t$] The set of grid cells that the predicted intruder may occupy at grid time $t$, the cells within the inflated radius plus half a cell of the predicted mean; the replanner treats them as blocked or expensive in time layer $t$ of the space-time graph.
\item[Reconnection] The return to the nominal path after an avoidance maneuver: a straight flight to a waypoint $\pi_i[j]$ ahead on the plan, arriving at grid time $t_0^i + j + \Delta$ with a delay $\Delta \le \Delta_{\max}$, admissible when it is slow enough and clear of obstacles, of the inflated prediction and of the teammates' intended positions.
\item[Conflict re-check] The run of the conflict detector over the padded paths of all drones, from the current grid time on, after any change of a path in force (a replan, a repair or a delayed reconnection).
\item[Partial replanning (local \cbs)] The repair of a conflict found by the re-check by running \cbs on the two agents in conflict only, from the next grid time, with the paths of all other drones reserved.
\item[Activation length $\ell_{\mathrm{act}}$] The link length, below the communication range $R_{\mathrm{comm}}$, at which the soft range term of the preferred velocity starts to pull a drone toward its nearest teammate.
